\chapter{Embeddings and Representation Learning}
\label{chap:embeddings}

\begin{tldr}
This chapter is the program's home for \emph{production embedding diagnostics}: how embedding spaces fail, and how you detect the failure before users do. It covers the collapse taxonomy (complete, dimensional, cluster) with SVD and spectral diagnostics, similarity-metric and pooling pitfalls, how tokenization shapes embedding quality, and the memory and quantization arithmetic for embedding tables and vector stores. Embedding \emph{theory}---the Word2Vec, GloVe and FastText derivations, the Levy--Goldberg matrix-factorization view, anisotropy, Matryoshka representation learning, embedding evaluation and embedder selection---is canonical in [NLP~3]. End-to-end retrieval \emph{system design}---funnel budgets, hybrid retrieval, reranking, index sizing---is canonical in [SR~4]. This chapter assumes both and concentrates on the failure surface between them.
\end{tldr}

\section{Embedding Fundamentals}
\label{sec:embedding_fundamentals}

\textbf{Embedding.} A learned mapping $f: \mathcal{X} \to \R^d$ from a discrete or high-dimensional space to a lower-dimensional continuous vector space where geometric relationships encode semantic similarity.

\textbf{Key property}: Similar items have similar embeddings.
\begin{equation}
\text{sim}(x_1, x_2) \approx \text{sim}(f(x_1), f(x_2))
\end{equation}

Why this is learnable at all---the distributional hypothesis, the count-based and PMI view of co-occurrence, and the objectives that turn them into trained vectors---is canonical in [NLP~3]. What this chapter builds on is the \emph{shape of the artifact}. The simplest embedding is a learned lookup table

\begin{equation}
\mE \in \R^{|V| \times d}, \quad \text{embed}(i) = \mE[i, :]
\end{equation}

with $|V| \times d$ parameters. That product is the unit of nearly every estimate in this chapter: an embedding table is memory before it is anything else, and a vector index is the same product with $|V|$ replaced by the corpus size and $d$ traded against precision.

%==============================================================================
\section{Static Embeddings in Production}
\label{sec:word_embeddings}
%==============================================================================

Word2Vec (skip-gram and CBOW, trained with negative sampling from a noise distribution $P_n(w) \propto f(w)^{0.75}$), GloVe (weighted least squares on log co-occurrence counts), and FastText (a word vector as the sum of its character $n$-gram vectors) are the three static embedding families. Their objectives, the Levy--Goldberg result that skip-gram with negative sampling implicitly factorizes a shifted PMI matrix, the linear analogy structure, and the polysemy limitation---one vector per word regardless of context---are derived in [NLP~3], and that is where to prepare the theory questions. This section keeps only what the production side of an interview asks: when a static table is still the right artifact.

\textbf{Static embeddings still win on four axes.} Contextual encoders win everywhere else, and the default for quality-sensitive work is a contextual or instruction-tuned embedder---but each of these is a real production constraint, not nostalgia:

\begin{itemize}
    \item \textbf{Latency.} A lookup is a memory access---microseconds---while a contextual encoder is a forward pass, milliseconds to tens of milliseconds. Under a sub-millisecond budget (autocomplete, per-request feature assembly) there is no contest.
    \item \textbf{Throughput and feature plumbing.} When a word or ID is one categorical feature among hundreds in a ranking model running at millions of events per second, an $O(1)$ table is the only affordable representation. This is why recommendation stacks are dominated by embedding tables rather than encoders (Chapter~\ref{chap:recommendation}).
    \item \textbf{Memory footprint.} A 50K-word, 100-dimensional FastText table is $\sim$20\,MB; a distilled BERT is $\sim$250\,MB. On-device deployment is often decided on this line alone.
    \item \textbf{OOV and morphology.} FastText composes unseen words from character $n$-grams, so a new morphological form gets a reasonable vector with no fine-tuning data---valuable for morphologically rich and low-resource languages where fine-tuning an encoder is not an option.
\end{itemize}

The pragmatic middle ground is \textbf{distillation into a table}: run the contextual encoder offline over the vocabulary or the item corpus, freeze the outputs, and serve them as a static lookup. That buys contextual quality at lookup cost for everything that can be scored offline---candidate generation, cached item vectors, batch scoring---while live encoding is spent only where context actually varies (the query side, and reranking). It is the same asymmetry retrieval systems exploit: corpus-side embedding is a batch job, query-side embedding is in the hot path ([SR~4]).

\begin{warningbox}
Static embeddings are not a modern \emph{default}. Reaching for Word2Vec when a pre-trained sentence encoder fits the budget is a red flag---and so is asserting that static embeddings are obsolete when a latency, memory, or QPS constraint plainly rules encoders out. The defensible answer names the constraint first and the artifact second. The quality argument for contextual embeddings (polysemy, word-sense disambiguation, sentence-level semantics) is [NLP~3].
\end{warningbox}

%==============================================================================
\section{Contextual Embeddings}
\label{sec:contextual}
%==============================================================================

ELMo (Peters et al., 2018) produced the first widely adopted contextualized vectors---a frozen two-layer biLSTM over a character CNN, mixed per task as $\text{ELMo}_k = \gamma \sum_{\ell} s_\ell \vh_{k,\ell}$---and established the layer-wise finding that lower layers carry syntax while upper layers carry semantics. That story, and the contextual-versus-static comparison it motivates, is canonical in [NLP~3]. What matters here is the structure of the input embedding that every encoder since BERT inherited, because it is what you count when you size a model and what you perturb when you debug one.

\subsection{BERT-style Input Embeddings}

BERT's input representation is the \emph{sum} of three separate embedding components:

\begin{equation}
\vx_i = \mE_{\text{token}}[w_i] + \mE_{\text{segment}}[\text{seg}_i] + \mE_{\text{position}}[i]
\end{equation}

\begin{itemize}
    \item \textbf{Token embeddings} ($\mE_{\text{token}} \in \R^{30{,}000 \times 768}$): WordPiece with a 30K vocabulary. Subword units handle OOV gracefully---``playing'' $\to$ ``play'' + ``\#\#ing''---and the ``\#\#'' prefix makes reconstruction unambiguous. The 30K size is the usual sweet spot: large enough that most English words are single tokens (shorter sequences), small enough to keep the table tractable.
    \item \textbf{Segment embeddings} ($\mE_{\text{segment}} \in \R^{2 \times 768}$): whether a token belongs to sentence A or B, which is what enables sentence-pair tasks (NLI, QA, paraphrase).
    \item \textbf{Position embeddings} ($\mE_{\text{position}} \in \R^{512 \times 768}$): learned, not sinusoidal, and capped at 512 tokens. This is a hard limit---longer sequences require architectural changes (Longformer, BigBird).
\end{itemize}

\begin{keyinsight}[Why BERT's Embedding Design Became the Template]
BERT established the pattern of \emph{summing} multiple embedding types (token + position + task-specific metadata) that nearly every subsequent encoder adopted. RoBERTa dropped segment embeddings (no NSP), ALBERT factorized the token embedding matrix, and DeBERTa disentangled content from position---but the ``sum of embeddings as input'' paradigm remained. When an interviewer asks about any encoder's input layer, start from this three-component template and explain what changed.
\end{keyinsight}

\textbf{Which layer should you take features from?} The last layer after fine-tuning; the second-to-last layer or the mean of the last four for frozen general-purpose features; not the raw \texttt{[CLS]} vector, which is a poor similarity representation without fine-tuning. The mechanism behind that last point---anisotropy, the narrow cone that raw LM features occupy---is canonical in [NLP~3], and it is also why the collapse thresholds in Section~\ref{sec:embedding_properties} must be calibrated to how the embeddings were trained rather than applied as absolutes.

%==============================================================================
\section{Sentence and Document Embeddings}
\label{sec:sentence_embeddings}
%==============================================================================

How sentence encoders are built and chosen---Sentence-BERT's siamese fine-tuning, SimCSE's dropout-augmented contrastive objective, the modern encoder and decoder-embedder families (E5, BGE, GTE and the 1--8B decoder embedders), Matryoshka representation learning, and the evaluation protocol (MTEB as a shortlist filter plus a judged domain set)---is canonical in [NLP~3]. The production residue that belongs here is \textbf{pooling}: the one part of the pipeline you own at serving time, and a frequent silent cause of bad vectors.

\begin{table}[H]
\centering
\caption{Sentence embedding pooling strategies}
\begin{tabularx}{\textwidth}{lXX}
\toprule
\textbf{Strategy} & \textbf{Method} & \textbf{Notes} \\
\midrule
{[CLS]} token & Use {[CLS]} representation & Designed for this, but needs fine-tuning \\
Mean pooling & Average all token embeddings & Simple, often effective \\
Max pooling & Element-wise max across tokens & Captures salient features \\
Attention pooling & Weighted average with learned weights & Task-adaptive \\
Last token & Use final token embedding & For autoregressive models \\
\bottomrule
\end{tabularx}
\end{table}

Two pooling failures recur in production, and both are invisible without an explicit check. First, \textbf{pooling mismatch}: encoding documents with mean pooling and queries with \texttt{[CLS]}, or normalizing on one side only, yields plausible-looking cosine values and quietly degraded ranking---the same failure class as embedding-version skew ([SR~4]). Second, \textbf{fragment dilution}: mean pooling averages over tokens, so text whose domain terms fragment into many subword pieces has its semantic signal diluted by those fragments (Section~\ref{sec:tokenization_embedding}). Pin the pooling strategy alongside the model version, and diff a sample of query and document vectors after every deployment.

%==============================================================================
\section{Embeddings from Other Modalities}
\label{sec:image_embeddings}
%==============================================================================

Image and multimodal embeddings come from different encoders but fail in the same ways, and every diagnostic in this chapter transfers unchanged.

\begin{itemize}
    \item \textbf{CNN penultimate layer} (ResNet-50 pool5, 2048-d; EfficientNet-B0, 1280-d): classification features reused as embeddings---cheap, but shaped by the label space they were trained on.
    \item \textbf{CLIP image tower} (512--1024-d, aligned with the text tower): enables zero-shot classification and cross-modal retrieval.
    \item \textbf{Self-supervised backbones} (SimCLR, DINO and DINOv2, MAE; 384--1536-d): the strongest label-free transfer, and DINOv2 features are a common frozen bank.
\end{itemize}

The methods themselves---contrastive, self-distillation and masked-prediction pre-training, and why the non-contrastive ones avoid collapse---are Chapter~\ref{chap:self_supervised}; cross-modal alignment is Chapter~\ref{chap:multimodal_learning}. They appear here only because collapse, dimensional under-use, and quantization loss are modality-agnostic: the SVD and uniformity diagnostics below read a DINOv2 feature bank exactly as they read a text index.

%==============================================================================
\section{Embedding Space Diagnostics}
\label{sec:embedding_properties}
%==============================================================================

\subsection{Similarity Metrics}

\begin{table}[H]
\centering
\caption{Embedding similarity metrics}
\begin{tabularx}{\textwidth}{lXX}
\toprule
\textbf{Metric} & \textbf{Formula} & \textbf{When to Use} \\
\midrule
Cosine & $\frac{\vu \cdot \vv}{\|\vu\| \|\vv\|}$ & Normalized embeddings, semantic similarity \\
Dot product & $\vu \cdot \vv$ & When magnitude matters, retrieval scores \\
Euclidean & $\|\vu - \vv\|_2$ & When embeddings are not normalized \\
\bottomrule
\end{tabularx}
\end{table}

If vectors are L2-normalized, cosine and dot product are the same ranking function, and dot product is cheaper for ANN search---so normalize at index time and score with dot product. The bug to watch for is a \textbf{metric mismatch}: training with cosine and serving with unnormalized dot product (or an index configured for L2 distance) changes the ranking without raising an error. Confirm the norms of stored vectors, not the config file.

\subsection{Embedding Collapse}

\textbf{Embedding Collapse.} When embeddings converge to a low-dimensional subspace or constant vector, losing representational capacity.

\textbf{Signs}:
\begin{itemize}
    \item All embeddings have high cosine similarity
    \item Singular values of embedding matrix decay rapidly
    \item Poor downstream performance despite low training loss
\end{itemize}

\textbf{Prevention}:
\begin{itemize}
    \item Contrastive loss (pushes negatives apart)
    \item Regularization (L2 on embeddings)
    \item Normalization (project to unit sphere)
    \item Decorrelation losses (Barlow Twins)
\end{itemize}

\textbf{Anisotropy caveat.} Uniformly high pairwise cosine is only evidence of collapse for embeddings trained with a contrastive or metric objective. Raw pre-trained LM features occupy a narrow cone and routinely average 0.5--0.9 cosine while still ranking neighbors usefully; the mechanism and the fixes are canonical in [NLP~3]. Judge such spaces by the \emph{gap} between intra-class and inter-class similarity, not by the absolute level.

\subsection{Linear Structure, and How Far to Trust It}

Well-trained embedding spaces carry approximate linear regularities:
\begin{equation}
\text{embed}(\text{``king''}) - \text{embed}(\text{``man''}) + \text{embed}(\text{``woman''}) \approx \text{embed}(\text{``queen''})
\end{equation}
and [NLP~3] derives why (the PMI-factorization view) and how the analogy task is actually scored. The caveat worth carrying into an interview: the canonical result depends on the evaluation protocol, which excludes the input words from the candidate set; without that exclusion the nearest neighbor of the arithmetic result is often ``king'' itself (Nissim et al., 2020). Present it as evidence of approximate linear structure, never as exact vector arithmetic.

%==============================================================================
\section{Tokenization and Embedding Interaction}
\label{sec:tokenization_embedding}
%==============================================================================

\subsection{Subword Tokenization}

Modern models use subword tokenization (BPE, WordPiece, SentencePiece) rather than word-level or character-level tokenization. The tokenizer algorithms themselves are covered in Chapter~\ref{chap:nlp_architectures}; what is canonical \emph{here} is how the choice propagates into embedding quality, because that is where it shows up as a production defect.

\textbf{BPE (Byte Pair Encoding)} iteratively merges the most frequent character pairs to build a vocabulary. A word like ``unhappiness'' might be tokenized as \texttt{[un, happi, ness]}.

\begin{table}[H]
\centering
\caption{Tokenization strategy and embedding quality trade-offs}
\begin{tabularx}{\textwidth}{lXX}
\toprule
\textbf{Strategy} & \textbf{Advantages} & \textbf{Disadvantages} \\
\midrule
Word-level & Clean semantics per token & OOV problem, huge vocabulary \\
Character-level & No OOV, tiny vocabulary & Long sequences, weak semantics per token \\
BPE/WordPiece & No OOV, moderate vocab & Frequent words get single token; rare words get fragmented \\
Byte-level (GPT-2+) & Universal across languages & Some tokens are meaningless byte sequences \\
\bottomrule
\end{tabularx}
\end{table}

\textbf{Vocabulary size trade-off}: Larger vocabularies mean more embedding parameters ($|V| \times d$) and more single-token representations (better per-token semantics), but also larger memory footprint and sparser training signal per token. Common sizes: 30,522 (BERT, WordPiece), 50,257 (GPT-2, byte-level BPE), 32K (LLaMA 2), 128K (LLaMA 3)---the trend is toward larger vocabularies as models scale.

\textbf{Key interaction}: Rare words get split into many subword tokens. Each subword has its own embedding, and the model must compose them through attention layers to reconstruct the word's meaning. This means embedding quality for rare words depends on the model's depth and attention capacity, not just the embedding table itself.

%==============================================================================
\section{Interview Questions}
\label{sec:embedding_questions}
%==============================================================================

%--- Q1: embedding space diagnosis ---

\begin{interviewq}
\textbf{Q: Your embedding space shows that unrelated items cluster together. Diagnose and fix.}

\badgeLsix{L6} \quad \badgetype{Debugging} \quad \badgefreq{\freqmed}

\stronganswer

Unrelated items clustering together in embedding space means the model has failed to learn discriminative representations. I would diagnose the root cause before applying fixes.

\textbf{Step 1: Characterize the failure.} Compute the average pairwise cosine similarity for random pairs---if it's $>0.8$ for a contrastively trained model, this is likely \textit{embedding collapse} (all embeddings converging to a similar region), not just local clustering errors. (Caveat: raw pre-trained LM features are anisotropic and can average 0.5--0.9 cosine without being collapsed, so calibrate the threshold to how the embeddings were trained.) Plot the singular value spectrum of the embedding matrix: if the effective rank is much lower than the embedding dimension (e.g., 95\% of variance captured by 10 out of 768 dimensions), the embeddings have collapsed to a low-dimensional subspace.

\textbf{Step 2: Check for collapse causes.} (a) \textit{Insufficient negative mining}: If contrastive training uses only random negatives, the model can achieve low loss by mapping everything to a general ``content'' region. Hard negatives force the model to make fine-grained distinctions. (b) \textit{Temperature too high}: In InfoNCE loss, $\loss = -\log \frac{\exp(\text{sim}^+/\tau)}{\sum \exp(\text{sim}/\tau)}$, a high $\tau$ makes the softmax ``softer,'' reducing the penalty for similar negative scores. Lower $\tau$ (e.g., 0.05--0.1) forces sharper discrimination. (c) \textit{Learning rate too high}: The model may have overshot to a degenerate minimum. (d) \textit{Batch size too small}: With in-batch negatives, small batches provide weak negative signal.

\textbf{Step 3: Check for data issues.} (a) \textit{Duplicate or near-duplicate positives}: If the training data has many near-identical pairs labeled as positive, the model learns to cluster everything together. (b) \textit{Label noise}: Incorrect positive labels (unrelated items paired together) directly teach the model to cluster unrelated items. Audit a random sample of training pairs. (c) \textit{Distributional imbalance}: If 80\% of items belong to one category, the model may learn a representation dominated by that category.

\textbf{Fixes, in order of effort:} (1) Add hard negatives (BM25 or model-mined). (2) Lower the temperature ($\tau$). (3) Increase batch size for more in-batch negatives. (4) Add a uniformity loss term (VICReg-style) that penalizes embeddings for being too similar. (5) Apply L2 normalization to the embedding output to project onto the unit hypersphere. (6) Clean the training data.

\redflags
\begin{itemize}
    \item ``Just increase the embedding dimension''---dimension doesn't help if the model collapses to a low-rank subspace.
    \item ``Retrain with more data''---this doesn't address the structural cause (collapse, bad negatives, label noise).
    \item Not distinguishing between full collapse (all items cluster) and partial failure (specific categories cluster wrongly).
\end{itemize}

\interviewtest{Systematic embedding debugging methodology---can you move from symptom (bad clusters) to diagnosis (collapse vs. data issue vs. training issue) to targeted fix, rather than guessing?}

\goodgreat
\begin{itemize}
    \item \badgeLfive{L5} Identifies embedding collapse as a likely cause, proposes hard negatives and temperature tuning.
    \item \badgeLsix{L6} Uses singular value analysis to characterize the collapse, distinguishes between global collapse and local clustering errors, and proposes uniformity-based regularization (VICReg, Barlow Twins).
    \item \badgeLseven{L7} Connects to the theory of representation learning---the uniformity-alignment framework (Wang \& Isola, 2020) where good embeddings need both alignment (similar items close) and uniformity (embeddings spread over the hypersphere). Proposes monitoring both metrics during training as early-warning signals.
\end{itemize}

\followups{``How do you detect collapse early during training?'' ``What is the uniformity-alignment trade-off?'' ``How does the VICReg loss prevent collapse without negatives?''}
\end{interviewq}

%--- Q2: embedding-side choices for a large index (consolidated; design is [SR 4]) ---

\begin{interviewq}
\textbf{Q: You are putting 100M item embeddings behind a vector index. Which embedding-side choices do you make---dimension, precision, refresh---and how would you know if one of them silently cost you recall?}

\badgeLsix{L6} \quad \badgetype{Estimation} \quad \badgefreq{\freqhigh}

\stronganswer

The end-to-end design---funnel stages, hybrid retrieval with BM25, reranking, latency budget, index family---is the canonical retrieval design question and belongs in that answer ([SR~4]); model \emph{selection} across the embedder landscape is [NLP~3]. What is specific to the embeddings is three decisions, each with an arithmetic and a detection story. The detection half is what candidates skip, and it is the half that matters: every one of these choices degrades quality \emph{silently}.

\textbf{Dimension is a multiplicative infrastructure cost.} Storage, index memory, and ANN latency all scale roughly linearly in $d$. At 100M vectors: 768 dims in FP32 is $100\text{M} \times 768 \times 4 = 307$\,GB; at 256 dims, 102\,GB. That is the difference between a sharded cluster and a single RAM-resident index, so dimension gets budgeted like latency, not chosen by default. Matryoshka-trained embeddings turn it into a post-training knob---truncate, renormalize, index the smallest dimension that passes your domain evaluation, and optionally re-score a shortlist at full dimension ([NLP~3] has the training details).

\textbf{Precision compounds with dimension.} FP16 halves FP32. Scalar int8 quantization gives another $2\times$ over FP16 ($100\text{M} \times 256 \times 1$\,byte $= 25.6$\,GB at 256 dims), plus a per-vector scale and zero-point---at two FP32 values per vector that is 0.8\,GB, negligible. Product quantization at 64 bytes per vector is 6.4\,GB, roughly $48\times$ below the FP32 768-dim baseline, with the full vectors kept on SSD for re-scoring. Compression and truncation multiply: choose one point on that grid, then measure it.

\textbf{Refresh and versioning.} New items embed on ingest and upsert into the index---at 1M new items/day this is $\sim$12 writes/second, operationally trivial. The expensive event is a \emph{model} change, which invalidates every stored vector: a v2 query vector scored against v1 document vectors returns plausible cosine values and garbage ranking, with no error anywhere. Stamp every vector with an \texttt{embedding\_version}, pin the query encoder to the index version as a serving invariant, and migrate through a shadow index with dual writes and a backfill before cutover (rollout mechanics: [SR~4]).

\textbf{How you detect the silent cost of each:}
\begin{enumerate}
    \item \textbf{Compression loss vs.\ ANN loss, separately.} Take 1--5K sampled queries, compute exact brute-force top-$k$ over the full-precision vectors as ground truth, then (a) exact search over the \emph{quantized} vectors to isolate quantization loss, and (b) the real index to isolate ANN approximation loss. Reporting one blended ``recall'' number conflates two failures with different fixes.
    \item \textbf{Truncation.} At each candidate dimension, measure rank correlation between full-$d$ and truncated-$d$ neighbor lists plus recall@$k$ on the judged domain set; take the knee, not the leaderboard default.
    \item \textbf{Spectral drift.} Track the singular-value profile and mean pairwise cosine of a rolling sample of newly embedded items against a reference epoch. A spectrum that flattens or concentrates signals that the item distribution or the embedder changed---the same dimensional-collapse diagnostic used at training time (Section~\ref{sec:embedding_properties}), reused as a production monitor.
    \item \textbf{Version skew.} Alert on any query whose encoder version does not match the index's stamped version. This is a two-line check that catches the most damaging and least visible failure in the list.
\end{enumerate}

\redflags
\begin{itemize}
    \item Quoting index sizes without the arithmetic---dimension, bytes per component, and vector count are all given, so the number is computable.
    \item Treating quantization recall loss and ANN recall loss as a single number.
    \item Re-embedding into a live index, or serving a new query encoder before the backfill completes.
    \item Answering the whole system design here---the end-to-end retrieval design is a separate question ([SR~4]); this one is scoped to the embedding artifact.
\end{itemize}

\interviewtest{Whether you can price the embedding artifact (dimension $\times$ precision $\times$ count) and, more importantly, whether you know that each compression and refresh decision needs its own measurement---because all of them fail without raising an error.}

\goodgreat
\begin{itemize}
    \item \badgeLfive{L5} Sizes the index correctly, picks a dimension and a quantization scheme, and mentions re-embedding when the model changes.
    \item \badgeLsix{L6} Adds MRL truncation as a post-training knob, separates quantization loss from ANN loss in the measurement plan, and makes version stamping plus query/index pinning an explicit invariant.
    \item \badgeLseven{L7} Treats migration economics as the dominant operational risk (re-embed cost, shadow index, validated cutover with the old index warm for rollback), and turns the training-time collapse diagnostics into continuous production monitors with alert thresholds.
\end{itemize}

\followups{``What breaks first if you truncate 768 to 128?'' ``How do you roll out a new embedding model with zero relevance downtime?'' ``Your index recall is 0.92---is that the embeddings or the ANN parameters?''}
\end{interviewq}

%--- Q3: memory arithmetic ---

\begin{interviewq}
\textbf{Q: Estimate the memory for an embedding table with 50K vocab, 768 dimensions in FP16 vs INT8.}

\badgeLfive{L5} \quad \badgetype{Estimation} \quad \badgefreq{\freqmed}

\stronganswer

\textbf{FP32 (baseline for reference):} $50{,}000 \times 768 \times 4$ bytes = $50{,}000 \times 3{,}072$ = 153,600,000 bytes $\approx$ \textbf{146.5 MB}.

\textbf{FP16:} $50{,}000 \times 768 \times 2$ bytes = 76,800,000 bytes $\approx$ \textbf{73.2 MB}. Half the FP32 size---this is the standard for most model serving.

\textbf{INT8:} $50{,}000 \times 768 \times 1$ byte = 38,400,000 bytes $\approx$ \textbf{36.6 MB}. Plus quantization parameters: for per-row quantization (one scale and zero-point per vocabulary entry), add $50{,}000 \times 2 \times 4$ bytes (scale + zero-point in FP32) = 400 KB---negligible overhead.

\textbf{Practical context:} A 50K vocab at 768-dim is small (BERT-sized). For modern LLMs with 128K vocab and 4096-dim embeddings: $128{,}000 \times 4{,}096 \times 2$ bytes (FP16) = 1.05 GB just for the embedding table. Whether to tie input and output embedding weights (saving one copy) depends on scale: tying is common in \textit{smaller} models where the table is a large fraction of total parameters (GPT-2, Gemma), while most large models (the LLaMA family, GPT-3-scale and beyond) leave them untied because the relative saving is small and untying adds flexibility. Either way, vocabulary size is a significant architectural decision.

\textbf{Memory beyond the embedding table:} In production, the embedding table is often a small fraction of total memory. The dominant costs are: (1) model weights (for a 7B model, $\sim$14 GB in FP16), (2) KV cache during inference ($\sim$1--4 GB for a single sequence at moderate context length; at production batch sizes the KV cache commonly reaches tens of GB and becomes the binding memory constraint for serving), (3) activation memory during training. The embedding table matters most when it's the primary artifact---e.g., in a pure retrieval system where you store embeddings for millions of items in a vector database.

\redflags
\begin{itemize}
    \item Getting the byte sizes wrong: FP32 = 4 bytes, FP16/BF16 = 2 bytes, INT8 = 1 byte.
    \item Forgetting quantization overhead (scale/zero-point per group or per row).
    \item Not contextualizing the estimate---73 MB is nothing for a GPU, but matters for edge deployment.
\end{itemize}

\interviewtest{Back-of-envelope calculation skills, understanding of precision formats, and ability to put numbers in practical context.}

\goodgreat
\begin{itemize}
    \item \badgeLfive{L5} Correct calculation for FP16 and INT8, with clear arithmetic.
    \item \badgeLsix{L6} Adds quantization parameter overhead, scales the calculation to modern LLM vocabulary sizes, and discusses weight tying.
    \item \badgeLseven{L7} Discusses practical deployment implications---how embedding table size interacts with the parallelism strategy (in Megatron-style tensor parallelism the table is \textit{sharded along the vocabulary dimension}---vocab-parallel embedding---whereas pure data parallelism replicates it on every GPU, as it does all weights), how vocabulary size affects throughput (larger softmax over vocabulary), and the trade-off between vocabulary size and sequence length for a fixed compute budget.
\end{itemize}

\followups{``How does the output embedding (LM head) relate to the input embedding?'' ``What is weight tying and why is it used?'' ``How would you quantize an embedding table without losing retrieval quality?''}
\end{interviewq}

%--- Q4: embedding collapse ---

\begin{interviewq}
\textbf{Q: What is embedding collapse and how do you detect/prevent it?}

\badgeLsix{L6} \quad \badgetype{Conceptual} \quad \badgefreq{\freqmed}

\stronganswer

\textbf{Definition.} Embedding collapse occurs when a learned embedding function maps all (or most) inputs to a small region of the embedding space, losing discriminative power. In the extreme case, all embeddings converge to a single point. In milder forms, embeddings occupy a low-dimensional subspace of the full embedding space (dimensional collapse).

\textbf{Why it happens.} The answer depends on the objective, and interviewers probe this distinction. For \textit{positive-only} objectives (e.g., BYOL or SimSiam with the predictor and stop-gradient removed), collapse \textit{is} a degenerate global minimum: mapping every input to the same vector makes every positive pair perfectly aligned at zero loss, and nothing in the loss opposes it---the model can reach this state by learning a constant function, which is easier than learning discriminative features. For contrastive losses \textit{with negatives} (InfoNCE), collapse is \textit{not} the global minimum: if all embeddings are identical, positive and negative similarities are equal, so the loss is $\log N$---chance level, strictly worse than what well-separated embeddings achieve. This repulsive term is precisely why contrastive losses resist collapse. When collapse does occur under a contrastive loss, it is an \textit{optimization} failure (negatives too easy to supply a repulsive gradient, temperature or learning rate poorly set), not a property of the loss's global structure.

\textbf{Types of collapse:}
\begin{itemize}
    \item \textbf{Complete collapse}: All embeddings map to (approximately) the same vector. Cosine similarity between random pairs is $>0.99$.
    \item \textbf{Dimensional collapse}: Embeddings use only a few dimensions of the available space. The embedding matrix has low effective rank. Detected by SVD: if the top $k \ll d$ singular values capture $>95\%$ of variance, the embeddings are effectively $k$-dimensional.
    \item \textbf{Cluster collapse}: Embeddings collapse into a few clusters regardless of true semantic categories. All items within each cluster are indistinguishable.
\end{itemize}

\textbf{Detection metrics:}
\begin{enumerate}
    \item \textbf{Average pairwise cosine similarity}: Sample 10K random pairs; for contrastively trained embeddings a healthy range is 0.0--0.3, and $>$0.7 signals collapse. \textit{Anisotropy caveat}: raw pre-trained LM features (no contrastive fine-tuning) occupy a narrow cone and routinely average 0.5--0.9 cosine while still ranking neighbors usefully---applied naively, this threshold would misdiagnose every vanilla-BERT feature space as collapsed. Use it only for embeddings trained with a contrastive or metric objective.
    \item \textbf{Singular value spectrum}: Compute SVD of the embedding matrix. Plot singular values on a log scale. A sharp drop-off (first few singular values dominate) indicates dimensional collapse.
    \item \textbf{Uniformity metric} (Wang \& Isola, 2020): $\mathcal{L}_{\text{uniform}} = \log \E[\exp(-2\|\vz_i - \vz_j\|^2)]$. More negative = more uniform distribution on the hypersphere = healthier embeddings.
    \item \textbf{Downstream task performance}: If training loss looks fine but downstream recall/accuracy doesn't improve, collapse is likely.
\end{enumerate}

\textbf{Prevention strategies:}
\begin{itemize}
    \item \textbf{Contrastive losses with sufficient negatives}: InfoNCE with large batch sizes and hard negatives creates strong repulsive forces.
    \item \textbf{Explicit decorrelation}: Barlow Twins loss penalizes off-diagonal elements of the cross-correlation matrix. VICReg adds explicit variance, invariance, and covariance regularization terms.
    \item \textbf{Asymmetric architectures}: BYOL and SimSiam avoid collapse through predictor networks and stop-gradient operations, even without negatives.
    \item \textbf{L2 normalization}: Projecting to the unit hypersphere prevents magnitude collapse. Combined with temperature scaling, this is standard practice.
    \item \textbf{Batch normalization on projections}: BN in the projection head (SimCLR) implicitly decorrelates dimensions, resisting dimensional collapse.
\end{itemize}

\redflags
\begin{itemize}
    \item ``Collapse means the model is underfitting''---collapse is a specific pathology of the embedding space geometry, not general underfitting.
    \item ``More training data prevents collapse''---collapse is a property of the loss function and architecture, not data quantity.
    \item Not distinguishing between complete collapse and dimensional collapse---they have different causes and fixes.
\end{itemize}

\interviewtest{Whether you understand the geometry of embedding spaces, can identify the degenerate solutions that loss functions admit, and know practical tools for detection and prevention.}

\goodgreat
\begin{itemize}
    \item \badgeLfive{L5} Defines collapse, explains why it's a problem, and proposes contrastive loss with hard negatives as prevention.
    \item \badgeLsix{L6} Distinguishes between complete and dimensional collapse, uses SVD for diagnosis, and knows about VICReg/Barlow Twins decorrelation losses.
    \item \badgeLseven{L7} Explains why BYOL/SimSiam avoid collapse without negatives (the predictor + stop-gradient creates an implicit negative force), connects to the uniformity-alignment framework, and discusses how to monitor collapse risk in production embedding training pipelines.
\end{itemize}

\followups{``Why doesn't BYOL collapse even without negative samples?'' ``What is the role of the projection head in preventing collapse?'' ``How does temperature in contrastive loss affect the collapse risk?''}
\end{interviewq}

%--- Q5: tokenization interaction ---

\begin{interviewq}
\textbf{Q: How does subword tokenization (BPE) interact with embedding quality?}

\badgeLfive{L5} \quad \badgetype{Conceptual} \quad \badgefreq{\freqhigh}

\stronganswer

Subword tokenization (BPE, WordPiece, SentencePiece) determines the granularity at which embeddings are learned, and this has direct consequences for embedding quality.

\textbf{Frequency-quality correlation.} Frequent words get single tokens (``the,'' ``is'') and therefore have dedicated embedding vectors trained on millions of examples. Rare words get split into multiple subword tokens (``photosynthesis'' $\to$ \texttt{[photo, synth, esis]}). Each subword has its own embedding, but the model must \textit{compose} them through subsequent layers (attention, FFN) to reconstruct the word's meaning. This composition is lossy---the model has seen fewer training examples of this specific combination, so the composed representation is noisier.

\textbf{Vocabulary size trade-off.} Larger vocabularies (128K in LLaMA 3 vs. 50K in GPT-2 or 32K in LLaMA 2) mean more words get single-token representations. Benefits: better per-token semantics, shorter sequences (fewer tokens per text $\to$ fits more context in the window, faster inference). Costs: larger embedding table (more parameters, more memory), each token seen less frequently during training (sparser signal). The optimal vocabulary size depends on the training data size---more data can support a larger vocabulary.

\textbf{Embedding quality impacts by token type:}
\begin{itemize}
    \item \textit{Common words} (single token): High-quality embeddings---many training examples, clear semantic signal.
    \item \textit{Subword fragments} (``un,'' ``ing,'' ``tion''): Embeddings capture morphological/syntactic patterns rather than standalone semantics. Quality depends on how consistently the fragment is used across contexts.
    \item \textit{Rare/domain-specific words} (multi-token): Quality depends on the composition mechanism. Transformers handle this better than RNNs because attention can directly connect all subword tokens.
    \item \textit{Non-English and code tokens}: Often heavily fragmented in English-centric tokenizers, leading to poor per-token embeddings and longer sequences. Multilingual tokenizers (SentencePiece trained on diverse data) mitigate this.
\end{itemize}

\textbf{Practical implications for embedding-based retrieval:} (1) If your domain has specialized vocabulary (medical terms, product names, code identifiers), the tokenizer may fragment these heavily, degrading retrieval quality. Fine-tuning on domain data helps because the model learns to compose these fragments correctly. (2) For sentence embeddings, the pooling strategy interacts with tokenization---mean pooling over many fragment tokens dilutes the semantic signal compared to pooling over fewer whole-word tokens.

\redflags
\begin{itemize}
    \item ``BPE is just about handling OOV words''---it fundamentally shapes the embedding space structure.
    \item ``Larger vocabulary is always better''---there's a data-dependent optimal size.
    \item Not knowing that BPE tokenization affects non-English languages disproportionately (more fragmentation $\to$ longer sequences $\to$ worse quality).
\end{itemize}

\interviewtest{Understanding of how the tokenization layer---often treated as a black box---directly impacts the learned representations and downstream task quality.}

\goodgreat
\begin{itemize}
    \item \badgeLfive{L5} Explains the frequency-quality correlation and the vocabulary size trade-off.
    \item \badgeLsix{L6} Discusses the composition mechanism (how Transformers reconstruct word meaning from fragments), the impact on retrieval quality, and multilingual tokenization issues.
    \item \badgeLseven{L7} Connects to recent work on byte-level models (MegaByte, byte-pair models) that bypass tokenization entirely, discusses how tokenizer-free approaches change the embedding quality landscape, and addresses the interaction between tokenization and compute efficiency (fewer tokens $\to$ cheaper training but potentially worse per-token embeddings).
\end{itemize}

\followups{``How would you choose vocabulary size for a new language model?'' ``How does tokenization affect multilingual models?'' ``What are Matryoshka embeddings and how do they address the dimension trade-off?''}
\end{interviewq}

%--- Q6: debugging embedding quality ---

\begin{interviewq}
\textbf{Q: How would you debug poor embedding quality?}

\badgeLfive{L5} \quad \badgetype{Debugging} \quad \badgefreq{\freqhigh}

\stronganswer

I would investigate in a systematic order, moving from cheap diagnostics to more expensive ones.

\textbf{Step 1: Visualize.} Use UMAP (preferred over t-SNE for preserving global structure) to project embeddings to 2D. Check if semantically similar items form clusters and if dissimilar items are separated. Color-code by known categories. If the visualization shows no meaningful structure, the embeddings are fundamentally broken.

\textbf{Step 2: Quantitative cluster analysis.} Compute average intra-class cosine similarity (should be high, e.g., $>0.6$) and inter-class cosine similarity (should be low, e.g., $<0.3$). If both are high, this is embedding collapse---though for raw un-fine-tuned LM features, uniformly high cosine is ordinary anisotropy rather than collapse, so judge by the intra-/inter-class \textit{gap} rather than the absolute values. If both are moderate, the embeddings are under-trained. If intra-class is low, the model hasn't learned category-level structure---check training data quality.

\textbf{Step 3: Nearest neighbor sanity check.} For 50--100 diverse query items, retrieve the top-10 nearest neighbors by cosine similarity. Manually inspect: do neighbors make semantic sense? Are there systematic errors (e.g., neighbors always share surface keywords but not deeper meaning)? This catches both embedding quality issues and similarity metric issues.

\textbf{Step 4: Singular value analysis.} Compute SVD of a sample embedding matrix. If the singular value spectrum drops off sharply (top 5\% of dimensions capture $>90\%$ of variance), the embeddings suffer from dimensional collapse---they aren't using the full representational capacity.

\textbf{Step 5: Downstream task evaluation.} Test embeddings on a simple proxy task (k-NN classification with frozen embeddings). If proxy performance is poor, the embeddings don't capture relevant features. If proxy performance is good but the actual downstream task is poor, the issue is in the downstream model, not the embeddings.

\textbf{Step 6: Training diagnostics.} Check loss curves (is loss still decreasing?), gradient norms (vanishing/exploding?), and learning rate schedule. If using contrastive learning, monitor the temperature parameter---too high ($>0.5$) makes the loss too easy, too low ($<0.01$) can cause instability.

\redflags
\begin{itemize}
    \item Jumping to ``retrain with more data'' before diagnosing the actual issue.
    \item Only using t-SNE visualization---this doesn't preserve global structure and can be misleading.
    \item Not computing quantitative metrics---``the UMAP looks fine'' is not a diagnosis.
\end{itemize}

\interviewtest{Systematic debugging methodology for embedding spaces, understanding of embedding geometry (cosine similarity, singular values, dimensionality), and the ability to distinguish between different failure modes.}

\goodgreat
\begin{itemize}
    \item \badgeLfive{L5} Proposes visualization + nearest neighbor check + downstream evaluation.
    \item \badgeLsix{L6} Adds quantitative metrics (intra/inter-class similarity, singular value analysis), distinguishes between collapse and underfitting, and checks training diagnostics.
    \item \badgeLseven{L7} Proposes a monitoring pipeline that runs these diagnostics automatically during training (not just post-hoc), and discusses how embedding quality metrics should be integrated into model release gates.
\end{itemize}

\followups{``What does embedding collapse look like in UMAP?'' ``How do you distinguish between bad embeddings and a bad evaluation metric?'' ``How would you set up automated embedding quality monitoring?''}
\end{interviewq}
